%%%%%%%%%%  scan idx 90  |  folio 77  %%%%%%%%%%
singulière. Dans ce deuxième cas, je n'ai pas fait la détermination complète, il manque un petit quelque chose. Avec les notations du petit papier, il faut prouver la formule (9) (mais comme \emph{égalité})~:
\begin{equation}\tag{9}
  f^*(i_*(\alpha))=j_*(\alpha C^{p-1}(F)),\qquad \alpha\in A(Y).
\end{equation}
On constate que les deux membres ont même image par \(f_*\), et même image par \(j^*\), de sorte que la question est équivalente à :
\begin{equation}\tag{9 bis}
  \Ker f_*\cap\Ker j^*=0.
\end{equation}
En tout cas, (9) est vraie en réduisant à \(G(K(X'))\) (ce qui détruit un peu de torsion, tout au plus), car elle résulte alors de (6)~; il n'y a donc guère de doute qu'elle ne soit vraie.

\dots Je te signale d'ailleurs que l'\emph{égalité} (9) résulte très facilement de (6) sans passer par l'ennuyeuse (6 bis)\dots

\begin{flushleft}
p.c.c.J-P.S.
\end{flushleft}

\newpage
